\section{Economic MPC and the design of the cost}
\label{sec:economic}

\subsection{Tracking and economic objectives}

The theory of the previous chapters was developed for \emph{regulation} and \emph{tracking}
problems: the control goals are first translated into reference values for the state or the
output, and the cost penalizes the distance from those references. Such a cost is positive
definite by construction, which is what the stability proof requires.

In many applications the goal is not a reference but a quantity with a direct meaning, usually
money: the electricity bill of a building, the cost of running a plant, the revenue of a service.
\emph{Economic MPC} writes that quantity directly as the cost, solves the problem, and applies
the first input. Nothing is tracked. The formulation is more natural, and the price is that the
assumptions of the theory are properties of the cost that an economic cost need not have: a bill
is not positive definite in the state, and the stability and optimality results of
\cref{sec:stability} do not apply as they stand. The case studies of \cref{sec:cases} are economic
MPC problems; their properties are studied, in the literature, with tools beyond this course.

\subsection{Several goals: the weighted sum}

A real controller has several goals at once: a bill to minimize, a storage to keep near a
reference, a service to deliver, a device to spare. The standard way to handle them is to write
one term per goal and to take a \emph{weighted sum},
\begin{equation}
  J = \alpha_1 J_1 + \alpha_2 J_2 + \dots + \alpha_r J_r, \qquad \alpha_j \ge 0 .
  \label{eq:weighted}
\end{equation}
Multiplying all the weights by the same positive constant does not change the minimizer, so
the weights can be normalized to sum to one, $\sum_j \alpha_j = 1$: only their ratios matter.

The set of cost pairs $(J_1, J_2)$ that admissible decisions can achieve, for two goals, is
sketched in \cref{fig:pareto}. Its lower-left boundary is the \emph{Pareto frontier}: the
solutions that cannot be improved in one goal without worsening the other. Minimizing the
weighted sum \eqref{eq:weighted} with a given pair of weights selects one point of the frontier
(this is the \emph{scalarization} of a vector optimization problem); sweeping the weights traces
the frontier --- all of it when the problem is convex, as in \cref{fig:pareto}, and only its
convex part otherwise. A rule-based controller, not being optimal, sits above the frontier; a better
algorithm moves the whole frontier towards the origin. Choosing the weights is therefore a
design decision about trade-offs, taken by whoever operates the process, not a mathematical one
--- and it is a decision that can be learned from data, which is one of the points where
reinforcement learning enters the picture.

\begin{figure}[t]
\centering
\begin{tikzpicture}[x=1cm,y=1cm,>=Stealth]
  \draw[->] (0,0) -- (6,0) node[right] {$J_1$};
  \draw[->] (0,0) -- (0,4.2) node[above] {$J_2$};
  \fill[mpc!10] plot[smooth] coordinates {(0.6,3.8) (1.0,2.3) (1.8,1.4) (3.0,0.9) (4.5,0.6) (5.5,0.55)} -- (5.5,4) -- (0.6,4) -- cycle;
  \draw[thick,mpc] plot[smooth] coordinates {(0.6,3.8) (1.0,2.3) (1.8,1.4) (3.0,0.9) (4.5,0.6) (5.5,0.55)};
  \node[mpc,font=\small,above right] at (3.0,0.95) {Pareto frontier};
  \fill[mpc] (0.6,3.8) circle (0.06) node[right,font=\small] {$(\alpha_1,\alpha_2)=(1,0)$};
  \fill[mpc] (5.5,0.55) circle (0.06) node[above,font=\small] {$(0,1)$};
  \fill[mpc] (1.8,1.4) circle (0.06) node[above right,font=\small] {$(0.5,0.5)$};
  \fill[sap] (3.4,2.6) circle (0.07) node[right,font=\small,sap] {rule-based controller};
  \draw[thick,dashed,ess] plot[smooth] coordinates {(0.4,3.0) (0.8,1.7) (1.6,0.95) (3.0,0.55) (5.5,0.3)};
  \node[ess,font=\small,anchor=north] at (2.3,0.45) {a better algorithm};
  \draw[->,ess] (2.3,0.47) -- (2.3,0.72);
\end{tikzpicture}
\caption{Two goals to minimize. The weighted sum with a given pair of weights selects one point
of the Pareto frontier; a rule-based controller lies above it; an algorithm with a larger
decision space pushes the frontier down.}
\label{fig:pareto}
\end{figure}

\subsection{The terminal value of a stored resource}

Economic problems have their own version of the terminal cost, and the building case study
(\cref{sec:building}) shows it. A controller that minimizes a bill over a window of one day
sells the energy stored in a battery before the window ends, because within the window that
sale is a profit and what happens afterwards does not exist. The behaviour is correct for the
problem written, and wrong for the application.

The remedy is to maximize the \emph{total value} over the window: the value of the state the
controller starts from, plus the profits minus the costs of what it plans, plus the value of
the state it leaves behind. The first term does not depend on the decisions, so maximizing the
total value is the same as minimizing
\begin{equation}
  J^{cost} - V(x_{\ik{N}}),
\end{equation}
where $V(x_{\ik{N}})$ is the price put on the resource still stored at the end of the window ---
for a battery, $V(x) = c^{avg}\,x$ with $c^{avg}$ the average purchase price of energy over the
day. This is the terminal cost of \cref{sec:optimality} seen from the side of the application:
it prices what happens after the horizon, and without it the controller is myopic in exactly
the way the theory predicts.

\begin{reading}
\refkey{Rawlings\_2017}, section 2.8, ``Economic Model Predictive Control'' (pp.~153--160): how
the economic problem differs from the tracking one, and which assumptions replace the positive
definite cost. \refkey{Boyd\_2004}, section 4.7, ``Vector optimization'' (pp.~174--187), in
particular 4.7.3 on Pareto optimal points and 4.7.4 on scalarization, which is the weighted sum
\eqref{eq:weighted} and the reason the weights can be normalized.
\end{reading}
